\section{Degradación de recursos: una sola ley para física, cómputo y red}
\label{sec:degradacion}

Los resultados anteriores acotan por separado tres cosas que parecen no tener
relación: cuánta fuerza entrega un contacto, cómo de bien está condicionado el
programa, y cuántas rondas necesita el acuerdo. Esta sección las reúne. La tesis
es que las tres degradan con tres escalares espectrales y que uno de ellos tiene
un umbral computable.

\begin{definicion}[Los tres recursos]
\label{def:recursos}
\[
  a=\sigma_{\min}\bigl(Q^{1/2}G_k\bigr),\qquad
  \delta=\operatorname{dist}\bigl(S_W\Wdem,\ \partial\mathcal W(\Ccal)\bigr),\qquad
  \eta=\text{tasa de mezcla del canal},
\]
respectivamente la \emph{autoridad} del acoplamiento esfuerzo--\emph{wrench},
el \emph{margen de factibilidad} de \cref{th:robusto} y la \emph{ergodicidad}
del grafo de información, ligada a $\lambda_2$ por \cref{th:mapa}.
\end{definicion}

\begin{teorema}[Degradación de recursos]
\label{th:degradacion}
Con $\varepsilon_\lambda>0$ la regularización de \eqref{eq:qp}:
\begin{enumerate}[label=(\roman*),leftmargin=2.2em,topsep=0.2em,itemsep=0.2em]
  \item \textbf{Umbral de parálisis.} Si $a<a^{\dagger}:=\sqrt{\varepsilon_\lambda}$,
  el residuo relativo de la dirección afectada excede $1/2$ y tiende a la unidad
  cuando $a\to0$: el programa, convexo y óptimo, renuncia a corregir.
  \item \textbf{Ley de esfuerzo.} Si $a>0$, entonces
  $\lVert\bm\mu^{\star}\rVert\le\lVert\tilde e\rVert/a$.
  \item \textbf{Ley de condicionamiento.} El número de condición del sistema KKT
  crece al menos como $1/\min\{a,\delta\}$.
  \item \textbf{Ley de rondas.} El número de rondas para alcanzar tolerancia
  $\epsilon$ satisface
  $R(\epsilon)\le\log(\epsilon_0/\epsilon)/\log(1/q)$ con factor de contracción
  compuesto $q\simeq1-c\min\{a^{2},\delta,\eta\}$.
\end{enumerate}
\end{teorema}

\begin{proof}
(i) es \cref{th:autoridad}(iii): en $a=\sqrt{\varepsilon_\lambda}$ el residuo
relativo $\varepsilon_\lambda/(a^{2}+\varepsilon_\lambda)$ vale exactamente
$1/2$, y la expresión es decreciente en $a$.
(ii) De $\mu_j^{\star}=\sigma_j\tilde e_j/(\sigma_j^{2}+\varepsilon_\lambda)$ se
sigue $|\mu_j^{\star}|\le\sigma_j|\tilde e_j|/\sigma_j^{2}=|\tilde e_j|/\sigma_j\le|\tilde e_j|/a$;
sumando cuadrados, $\lVert\bm\mu^\star\rVert\le\lVert\tilde e\rVert/a$.
(iii) es el transporte de la relación entre condicionamiento y distancia a la
infactibilidad: exacta en forma, con constante dependiente de la
instancia.
(iv) compone las tasas de las tres capas: la del reparto, gobernada por el
menor valor singular al cuadrado; la del mercado, por el margen; y la del
consenso, por la mezcla. La contracción del ciclo es la peor de las tres.
\end{proof}

\subsection{Estatuto de cada inciso}

No conviene presentar los cuatro incisos con la misma fuerza.

\begin{center}\small
\begin{tabular}{@{}lll@{}}
\toprule
Inciso & Estatuto & Base \\
\midrule
(i) umbral & demostrado exacto & \cref{th:autoridad}, régimen interior \\
(ii) esfuerzo & demostrado exacto & desigualdad directa \\
(iii) condicionamiento & exacto en forma, constante por instancia & transporte \\
(iv) rondas & ley de orden, constantes por calibrar & composición \\
\bottomrule
\end{tabular}
\end{center}

\begin{observacion}[Comprobación numérica]
\label{obs:degradacion-num}
La ley de esfuerzo se verificó sobre instancias aleatorias de cuatro direcciones
singulares: en todas, $\lVert\bm\mu^\star\rVert\le\lVert\tilde e\rVert/a$ con
holgura amplia, como corresponde a una cota que satura solo cuando una única
dirección domina. Para $\varepsilon_\lambda=0{,}04$, el residuo relativo pasa de
$0{,}038$ en $a=1$ a $0{,}20$ en $a=0{,}4$, exactamente $0{,}50$ en
$a=a^\dagger=0{,}2$ y $0{,}99$ en $a=0{,}02$. En la ley de rondas, el recurso
que limita cambia con el régimen: con $(a,\delta,\eta)=(0{,}3,\,0{,}8,\,0{,}7)$
limita la autoridad; con $(0{,}9,\,0{,}05,\,0{,}7)$, el margen; y con
$(0{,}9,\,0{,}8,\,0{,}02)$, el canal.
\end{observacion}

\subsection{Qué unifica}

\Cref{th:degradacion} da una sola explicación a fenómenos que en un banco de
pruebas aparecen como incidencias inconexas:

\begin{itemize}[leftmargin=1.4em,topsep=0.2em,itemsep=0.25em]
  \item una coalición que pierde contacto y cuyo controlador «no reacciona»
  —es $a\to0$ y la parálisis óptima de (i), no un fallo del solver;
  \item una configuración cuya solución exige esfuerzos desmesurados
  —es la ley (ii) con $a$ pequeño;
  \item un programa que converge mal cerca de la saturación
  —es (iii) con $\delta\to0$;
  \item un acuerdo que tarda cuando la flota se dispersa
  —es (iv) con $\eta$ pobre.
\end{itemize}

\begin{corolario}[Guarda espectral acoplada]
\label{co:guarda-global}
La condición de admisión $a\ge\kappa\sqrt{\varepsilon_\lambda}$ de
\cref{co:guarda} no es una caja calibrada: por \cref{th:degradacion} es
simultáneamente una cota del residuo, una cota del esfuerzo y una cota del
condicionamiento. Fijar $\kappa$ fija las tres.
\end{corolario}

\begin{observacion}[Los tres recursos son de naturaleza distinta]
\label{obs:recursos-distintos}
La unificación no debe ocultar una diferencia importante. $a$ es geométrico:
depende de dónde estén los contactos y hacia dónde miren los AMR, y la coalición
lo controla eligiendo puestos. $\delta$ es de demanda: depende de cuánto se pide
frente a lo que la coalición puede dar, y se controla frenando o reclutando.
$\eta$ es de red: depende del grafo, y se controla con la topología de
\cref{sec:topologia}. Que las tres degradaciones compartan forma no significa
que compartan remedio.
\end{observacion}
